\section{Modelado y discretización}
\subsection{Modelo continuo}
A partir del diagrama analógico de la guía oficial, con definición explícita
de los nodos que constituyen el vector de estados, se construirá
\begin{align}
 \dot{\mathbf{x}}(t) &= F\mathbf{x}(t)+G u(t),\\
 y(t)&=H\mathbf{x}(t)+J u(t).
\end{align}
Para una planta de segundo orden y una sola salida medida,
\(F\in\mathbb R^{2\times2}\), \(G\in\mathbb R^{2\times1}\),
\(H\in\mathbb R^{1\times2}\) y \(J\in\mathbb R\).
Las matrices numéricas, los valores R--C y las unidades se incorporarán
cuando se confirme el circuito físico; no se copiarán de un informe ajeno.

\subsection{Selección justificada del muestreo}
La selección de \(T_s\) precede al cálculo de \(A\) y \(B\). Se estudiarán la
frecuencia característica más rápida, el ancho de banda de diseño, la
relación de muestreo exigida en clase y las cuentas enteras del temporizador
\cite{samanta2024,fadali2020}. La frecuencia efectiva será
\begin{equation}
 f_s = \frac{1}{T_s},\qquad \omega_s=\frac{2\pi}{T_s}.
\end{equation}
El valor definitivo de \(T_s\) queda \emph{pendiente de verificación};
no se adopta por defecto el de los laboratorios anteriores.

\subsection{Modelo discreto con retenedor de orden cero}
Para mando constante entre muestras (ZOH),
\begin{align}
 A &= e^{FT_s}, & B &= \int_{0}^{T_s} e^{F\tau}G\,d\tau,\\
 C &= H, & D &= J. 
\end{align}
De ese modo,
\begin{equation}\label{eq:planta_discreta}
\mathbf{x}[k+1]=A\mathbf{x}[k]+Bu[k],\quad y[k]=C\mathbf{x}[k]+Du[k].
\end{equation}
El cálculo se verificará con \texttt{c2d(ss(F,G,H,J),Ts,'zoh')}.

\subsection{Controlabilidad y observabilidad}
Antes de asignar polos se comprobarán
\begin{equation}
 \mathcal C=[B\;AB],\qquad
 \mathcal O=\begin{bmatrix} C\\CA\end{bmatrix},
\end{equation}
y se exigirá \(\operatorname{rango}(\mathcal C)=
\operatorname{rango}(\mathcal O)=2\). La observabilidad debe verificarse
\emph{desde el único canal de salida} disponible en el experimento.
